\documentclass[12pt, letterpaper]{article}

\usepackage[margin=1in]{geometry}
\usepackage{titlesec}
\usepackage{enumitem}
\usepackage{xcolor}
\usepackage{mdframed}
\usepackage{parskip}
\usepackage{fancyhdr}
\usepackage{hyperref}
\usepackage{booktabs}
\usepackage{array}
\usepackage{amsmath}

% ── Colors ────────────────────────────────────────────────────────────────────
\definecolor{slideblue}{RGB}{52, 73, 94}
\definecolor{notebg}{RGB}{245, 247, 250}
\definecolor{accent}{RGB}{41, 128, 185}
\definecolor{slidenumber}{RGB}{150, 150, 150}

% ── Header / Footer ───────────────────────────────────────────────────────────
\pagestyle{fancy}
\fancyhf{}
\rhead{\textcolor{slidenumber}{\small DATS6401 -- FAERS Q4 2025}}
\lhead{\textcolor{slidenumber}{\small Kirtan Patel -- G34079442}}
\cfoot{\textcolor{slidenumber}{\thepage}}
\renewcommand{\headrulewidth}{0.4pt}

% ── Section styling ───────────────────────────────────────────────────────────
\titleformat{\section}
  {\large\bfseries\color{slideblue}}
  {Slide \thesection:}{0.5em}{}
  [\vspace{-4pt}\textcolor{accent}{\rule{\linewidth}{1.5pt}}\vspace{4pt}]

\titleformat{\subsection}
  {\normalsize\bfseries\color{accent}}
  {}{0em}{}

% ── Note box ──────────────────────────────────────────────────────────────────
\newmdenv[
  backgroundcolor=notebg,
  linecolor=accent,
  linewidth=1.5pt,
  leftline=true,
  rightline=false,
  topline=false,
  bottomline=false,
  innerleftmargin=12pt,
  innerrightmargin=8pt,
  innertopmargin=8pt,
  innerbottommargin=8pt,
  skipabove=6pt,
  skipbelow=6pt,
]{notebox}

% ─────────────────────────────────────────────────────────────────────────────
\begin{document}

\begin{center}
  {\LARGE\bfseries\color{slideblue} Presentation Speaker Notes}\\[6pt]
  {\large FDA Adverse Event Reporting System (FAERS) Q4 2025}\\[4pt]
  {\normalsize Kirtan Patel \quad G34079442 \quad George Washington University}\\[2pt]
  {\normalsize DATS6401 -- Visualization of Complex Data \quad Spring 2026}
\end{center}

\vspace{4pt}
\textcolor{accent}{\rule{\linewidth}{1pt}}
\vspace{8pt}

\begin{notebox}
\textbf{How to use this document:} Each section corresponds to one slide in the
presentation, in order. The speaker notes contain everything to say aloud,
including content omitted from the visual slide. Time estimates are provided as
a guide for the 20-minute presentation.
\end{notebox}

\vspace{10pt}

% =============================================================================
\section{Title Slide \hfill \textcolor{slidenumber}{\small $\sim$30 sec}}
% =============================================================================

\begin{notebox}
\textit{Say as you appear on screen -- keep it short.}
\end{notebox}

Good morning / good afternoon everyone. My name is Kirtan Patel, student ID
G34079442, and today I will be presenting my final term project for DATS6401 --
Visualization of Complex Data at George Washington University. The project is
titled \textbf{FDA Adverse Event Reporting System -- FAERS Q4 2025}, and covers
the complete pipeline from raw data ingestion all the way through an interactive
web-based dashboard deployed on Google Cloud Platform. Let me jump straight in.

% =============================================================================
\section{What is FAERS? \hfill \textcolor{slidenumber}{\small $\sim$1 min}}
% =============================================================================

The FDA Adverse Event Reporting System, commonly known as FAERS, is the United
States government's primary post-market drug safety surveillance database. It is
maintained by the FDA's Center for Drug Evaluation and Research, CDER, and
specifically by the Office of Surveillance and Epidemiology.

FAERS collects Individual Case Safety Reports -- or ICSRs -- submitted
voluntarily or mandatorily by healthcare professionals, consumers, and
pharmaceutical manufacturers. These reports document adverse drug reactions,
medication errors, and product quality complaints for all approved human drugs
and biological products sold in the United States.

What makes FAERS particularly important from a real-world standpoint is that it
has directly driven major regulatory actions. Analyses of FAERS data led to
black box warnings on several blockbuster drugs, the market withdrawal of
Vioxx and Bextra, and the creation of Risk Evaluation and Mitigation Strategies,
known as REMS programs. So this is not just academic data -- it has real
consequences for patient safety and regulatory policy.

In terms of scale, the cumulative database now contains over 20 million reports,
with more than 2 million new reports added annually. It is released every
quarter in both XML and ASCII flat-file formats and is freely downloadable
without registration from the FDA's public data portal.

% =============================================================================
\section{Dataset Scope and Structure \hfill \textcolor{slidenumber}{\small $\sim$1 min}}
% =============================================================================

Because the full cumulative FAERS database contains over 20 million reports and
would be computationally prohibitive to process on a single machine, this
project focuses on the most recent quarterly extract -- \textbf{Q4 2025},
covering October through December 2025. This extract contains approximately
200,000 to 400,000 individual safety reports, which is large enough for
statistically meaningful analysis while remaining manageable.

The data follows the international \textbf{ICH E2B v2.1 reporting standard},
which defines over 60 XML fields organized in a strict hierarchical structure.
After parsing and flattening this hierarchy, the data resolves into three core
relational tables:

\begin{itemize}[leftmargin=1.5em]
  \item \textbf{Demographics table} -- one row per unique safety report,
    containing patient age, sex, weight, country of occurrence, reporter type,
    all seriousness flags, and temporal fields like receive date and report date.
  \item \textbf{Drug table} -- multiple rows per report, since a single report
    can involve many drugs. Each row captures the drug name, its role
    (Suspect, Concomitant, Interacting, or Not Administered), route of
    administration, structured dosage fields, action taken with the drug, and
    the indication it was prescribed for.
  \item \textbf{Reaction table} -- multiple rows per report as well, since one
    patient can experience many adverse reactions. Each row holds a MedDRA
    Preferred Term describing the reaction and the outcome at last observation
    (Recovered, Fatal, Unknown, etc.).
\end{itemize}

All three tables are linked by the \texttt{safetyreportid} field, giving us a
one-to-many relationship: one report can have many drugs and many reactions.
After merging and processing, our final analytical dataset contains approximately
\textbf{385,000 unique reports}, \textbf{1.9 million drug records}, and
\textbf{1.3 million reaction records}.

% =============================================================================
\section{Features \hfill \textcolor{slidenumber}{\small $\sim$1 min}}
% =============================================================================

The dataset contains two broad categories of features. On the numerical side we
have seven continuous or discrete variables. \textbf{patientonsetage} captures
the patient's age at the time of the adverse event. \textbf{patientweight}
records weight in kilograms. \textbf{drugtreatmentduration} records how long the
drug was administered. The three dosage fields -- \textbf{drugseparatedosagenumb},
\textbf{drugintervaldosageunitnumb}, and \textbf{drugstructuredosagenumb} --
capture the dosing schedule, and \textbf{drugcumulativedosagenumb} records total
cumulative dose administered up to the first reaction.

On the categorical side we have over 36 fields, of which the most analytically
important are: \textbf{reporttype} (Spontaneous, Study, Other, Unknown),
\textbf{patientagegroup} (Neonate through Elderly), \textbf{qualification}
(reporter type -- Physician, Pharmacist, Other Healthcare Professional, Lawyer,
Consumer), \textbf{actiondrug} (what was done with the drug -- Withdrawn,
Dose Reduced, Not Changed, etc.), \textbf{drugindication} (the MedDRA term for
why the drug was prescribed), \textbf{occurcountry} (country where the event
happened), and \textbf{drugcharacterization} (Suspect, Concomitant,
Interacting, or Not Administered).

The \textbf{dependent variable} for analytical purposes is
\textbf{reactionoutcome}, which captures whether a patient recovered, did not
recover, recovered with lasting effects (sequelae), or experienced a fatal
outcome. We also defined a composite binary seriousness flag as a key grouping
variable throughout the analysis.

% =============================================================================
\section{Data Cleaning \hfill \textcolor{slidenumber}{\small $\sim$1.5 min}}
% =============================================================================

Raw FAERS data requires substantial preprocessing before it is usable for
analysis. The cleaning pipeline I implemented handles the following steps in
sequence.

First, the raw XML files are parsed using Python's event-based
\texttt{iterparse} method, which reads the XML tree incrementally without
loading the entire file into memory. For quarters where ASCII CSV files are
available, those are loaded as a fallback.

Second, all numeric codes are mapped to human-readable labels. For example,
sex code 1 becomes Male, code 2 becomes Female. Reaction outcome code 5 becomes
Fatal, code 1 becomes Recovered. Reporter qualification code 1 becomes
Physician, code 2 becomes Pharmacist, and code 5 becomes Consumer. The full
mapping covers every coded categorical field in the dataset.

Third, follow-up duplicate reports are removed. FAERS allows manufacturers to
submit multiple versions of the same case as new information becomes available.
We keep only the latest \texttt{safetyreportversion} per unique \texttt{caseid}
to avoid counting the same event multiple times.

Fourth, patient age is standardized to years. The raw field uses five different
unit codes -- Decade, Month, Week, Day, and Hour -- specified in a companion
field called \texttt{patientonsetageunit}. All values are converted to years
using the appropriate conversion factors. Ages are then capped between 0 and 120
years to remove physiologically implausible entries.

Fifth, patient weight is capped between 3 and 300 kilograms to remove obvious
data entry errors. Missing values for both age and weight are filled using a
\textbf{stratified median} approach -- the median is computed within subgroups
defined by age group and sex, so imputed values are contextually appropriate
rather than globally averaged.

Finally, the drug name field \texttt{medicinalproduct} is normalized to
uppercase and whitespace is trimmed to consolidate variant spellings of the same
drug. Where possible, brand names are mapped to their active ingredient names
using the \texttt{activesubstancename} field.

% =============================================================================
\section{Feature Engineering \hfill \textcolor{slidenumber}{\small $\sim$1 min}}
% =============================================================================

Beyond cleaning, several new features were engineered from the raw data.

The most important is the \textbf{Seriousness Score}, an integer from 0 to 6
computed by summing six binary flags: death, life-threatening event,
hospitalization, disabling condition, congenital anomaly, and other serious
criterion. This converts six separate yes/no columns into a single ordinal
measure of overall severity that can be used in regression, correlation, and
visualization.

Two derived binary flags -- \textbf{is\_serious} and \textbf{is\_fatal} -- were
created for filtering and grouping across all dashboard tabs.

Since each report can reference multiple drugs and multiple reactions,
\textbf{num\_drugs} and \textbf{num\_reactions} were engineered per report by
merging the drug and reaction tables back to the demographics table. These
become new numerical features capturing case complexity.

From the receive date field, temporal features were extracted including
\textbf{report\_month}, \textbf{day\_of\_week}, and a parsed datetime column
\textbf{receive\_dt} used for trend analysis.

For high-cardinality categorical fields like \texttt{medicinalproduct} which
contains over 10,000 unique values, and \texttt{drugadministrationroute},
\textbf{frequency encoding} was applied to make these fields analytically
tractable. The 27,000-plus MedDRA reaction terms were grouped into 8
\textbf{System Organ Classes} -- General Disorders, Gastrointestinal, Nervous
System, Respiratory, Skin, Musculoskeletal, Infections, and Product Issues --
for the Reaction Explorer tab. Finally, all seriousness flags were converted
from their original 1/2 binary coding to a standard 0/1 encoding.

% =============================================================================
\section{Outlier Detection \& Removal \hfill \textcolor{slidenumber}{\small $\sim$1 min}}
% =============================================================================

Outlier detection was performed using three complementary methods across all
seven numerical features.

The \textbf{IQR method} computes the interquartile range for each feature.
The lower fence is $Q1 - 1.5 \times IQR$ and the upper fence is
$Q3 + 1.5 \times IQR$. Any record falling outside these bounds is flagged as an
outlier. This method is robust to non-Gaussian distributions and is the
standard approach recommended for skewed medical data.

The \textbf{Z-Score method} standardizes each value using
$z = (x - \mu) / \sigma$. Records with an absolute z-score greater than 3,
meaning more than three standard deviations from the mean, are flagged as
extreme outliers. This method is sensitive to the underlying distribution and
works best in combination with the IQR approach.

\textbf{Hard capping} was applied as a preprocessing step for domain-specific
boundaries: age is capped at 0 to 120 years and weight at 3 to 300 kilograms.
These thresholds are physiologically justified -- any value outside them is
almost certainly a data entry error rather than a legitimate measurement.

Both the IQR and Z-Score methods are implemented interactively in the dashboard
Analytics tab. Users can select any numerical variable from a dropdown and
immediately see which records are flagged as outliers, with the count and
percentage of outliers annotated directly on the chart.

% =============================================================================
\section{Outlier Detection -- Dashboard Graphs \hfill \textcolor{slidenumber}{\small $\sim$30 sec}}
% =============================================================================

\begin{notebox}
\textit{This is a screenshot slide. Point to the visuals as you speak.}
\end{notebox}

On the left you can see the \textbf{IQR box plot} for patient age. The
annotation at the top shows IQR equals 16 years, Q1 equals 50, Q3 equals 66,
and a total of 34,622 records -- or 8.99 percent of the dataset -- fall outside
the 1.5 times IQR fence. These are shown as red cross markers beyond the
whiskers.

On the right is the \textbf{Z-Score scatter plot}. Each dot represents one
patient record. The x-axis shows the raw age value and the y-axis shows the
computed z-score. The two horizontal dashed red lines mark the plus and minus 3
sigma boundaries. The red circles near age zero represent the only extreme
outliers -- very young ages with z-scores below $-3$. The fact that almost all
points fall within the band confirms that hard capping during preprocessing
already removed the most extreme values before this analysis ran.

Both of these charts update in real time when you switch the variable dropdown
in the Analytics tab.

% =============================================================================
\section{Normality Testing \hfill \textcolor{slidenumber}{\small $\sim$1 min}}
% =============================================================================

To determine whether our numerical features follow a Gaussian distribution, I
applied the \textbf{Shapiro-Wilk test} to all seven numerical features.

The null hypothesis $H_0$ is that the data comes from a normal distribution.
We reject $H_0$ when the p-value is below 0.05. The Shapiro-Wilk test is
preferred over alternatives like the Kolmogorov-Smirnov test when the sample
size is moderate, and it remains valid even for large samples when the departure
from normality is substantial.

The results are unambiguous: all seven features reject normality at any standard
significance level. The dosage fields -- particularly
\texttt{drugstructuredosagenumb} and \texttt{drugcumulativedosagenumb} -- are
the most extreme cases. These fields span several orders of magnitude because
some drugs are dosed in micrograms and others in grams, and the distributions
have very long right tails. \texttt{patientonsetage} is the closest to symmetric
but still rejects normality, primarily due to the heavy concentration of
reports in the 50 to 70 age range.

These results directly motivated the data transformation step, where log
transformation and Min-Max normalization were applied to address the
non-Gaussian distributions before further statistical analysis.

In the dashboard, the \textbf{Q-Q Plot} for each variable shows sample quantiles
plotted against theoretical normal quantiles, with the Shapiro-Wilk p-value
annotated directly on the chart. If the data were normally distributed, all
points would fall on the diagonal red dashed line.

% =============================================================================
\section{Normality Testing -- Dashboard Graphs \hfill \textcolor{slidenumber}{\small $\sim$30 sec}}
% =============================================================================

\begin{notebox}
\textit{This is a screenshot slide. Point to both panels.}
\end{notebox}

On the left is the \textbf{Q-Q Plot} for patient age. You can see the blue
sample quantile curve deviates significantly from the red diagonal fit line,
particularly at both tails. The annotation confirms the Shapiro-Wilk test
result: Non-Normal with $p = 8.40 \times 10^{-50}$ -- a p-value so small it
is effectively zero. This definitively confirms the distribution is not Gaussian.

On the right is the \textbf{Normalization Comparison} chart -- one of the
interactive features I built. The blue bars show the raw age distribution and
the orange curve overlaid on the secondary y-axis shows how the Log$(x+1)$
transformation maps each raw value to a normalized value. You can switch this
dropdown between Min-Max, Z-Score, and Log to see how each method reshapes the
data. The left y-axis shows raw density in blue and the right y-axis shows the
normalized value in the matching color of the selected method.

% =============================================================================
\section{Dimensionality Reduction: PCA \hfill \textcolor{slidenumber}{\small $\sim$1 min}}
% =============================================================================

With seven numerical features confirmed as non-Gaussian and partially
correlated, I performed \textbf{Principal Component Analysis} to reduce
dimensionality and understand the underlying structure.

The full pipeline is: Z-Score standardize all seven features, compute the
covariance matrix, perform eigen decomposition, then select the minimum number
of principal components that together explain at least 90 percent of the total
variance.

The key finding is that the first two principal components explain the majority
of variance, but the cumulative curve rises gradually rather than sharply. This
is reflected in a high condition number, indicating that no single underlying
factor dominates the dataset. The seven features are largely independent
clinical measurements -- patient age, patient weight, and dosage fields capture
very different aspects of a case and do not collapse into a single latent
dimension. In the biplot, dosage-related features cluster in a different
direction from age and weight features, confirming their structural independence.

Three additional feature selection methods were applied to supplement PCA:
\textbf{Mutual Information Scoring} ranked categorical features by their
statistical relevance to the seriousness outcome.
\textbf{Variance Threshold Filtering} dropped any near-zero-variance columns
that add no discriminative information.
\textbf{Correlation Filtering} removed any feature pairs with an absolute
Pearson $r$ greater than 0.90 to reduce multicollinearity before modeling.

% =============================================================================
\section{Pearson Correlation Matrix -- Dashboard Graph \hfill \textcolor{slidenumber}{\small $\sim$1 min}}
% =============================================================================

\begin{notebox}
\textit{This is a screenshot slide. Walk through the key cells.}
\end{notebox}

This is the interactive \textbf{Pearson Correlation Matrix} from the dashboard,
now built in Plotly so you can hover any cell to read the exact correlation
coefficient. The matrix covers all 12 numerical variables including the seven
raw features plus the engineered variables: num\_drugs, num\_reactions,
seriousness\_score, and the individual flag fields.

Let me highlight the key observations:

The strongest correlation in the entire matrix is \textbf{Age Years versus
Patientonsetage} at $r = 0.94$. This is expected -- these two fields both
measure patient age, just computed through different pathways, which confirms our
age standardization pipeline worked correctly.

\textbf{Num Drugs versus Num Reactions} shows $r = 0.49$ -- a meaningful
positive correlation. Patients who are on more medications at the time of the
event also tend to have more documented adverse reactions, which aligns with
the pharmacological concept of polypharmacy.

\textbf{Seriousness Score versus Hospitalization Flag} shows $r = 0.62$, the
strongest relationship between the composite score and any individual flag,
which makes clinical sense -- hospitalization is the most common serious outcome
in the dataset.

Importantly, the dosage fields show near-zero correlations with patient
demographics, which explains why PCA yielded a high condition number -- these
feature groups are orthogonal to each other in the multivariate space.

% =============================================================================
\section{Conclusion \hfill \textcolor{slidenumber}{\small $\sim$1 min}}
% =============================================================================

To summarize the key findings from this project:

FAERS Q4 2025 contains over 385,000 unique safety reports. Of these,
\textbf{56.67 percent} are classified as serious adverse events and
\textbf{6.82 percent} resulted in a fatal outcome. Elderly patients consistently
show the highest seriousness scores and fatality rates across every analytical
view in the dashboard.

All seven numerical features were confirmed non-Gaussian by the Shapiro-Wilk
test, which directly justified applying log transformation and Box-Cox
transformation before any parametric analysis. The Normalization Comparison
chart in the dashboard makes this transformation visible to any user
interactively.

PCA confirmed that the seven features are largely independent clinical
measurements with no dominant latent factor -- the condition number is high and
variance is spread across multiple components. The positive correlation between
\texttt{num\_drugs} and \texttt{num\_reactions} demonstrates that polypharmacy
cases generate more complex adverse event profiles.

The Safety Signals tab implements a composite signal scoring formula weighting
report volume at 40 percent, serious rate at 35 percent, and fatal rate at 25
percent -- identifying the highest-risk drug--reaction pairs in Q4 2025 beyond
simple frequency counts.

Overall, the interactive dashboard makes these pharmacovigilance patterns
accessible to both clinical researchers and non-technical users through
real-time filtering, contextual tooltips, and clearly labeled charts across
10 analytical tabs.

For future work, I would integrate the full historical FAERS database across all
quarters, add formal disproportionality measures such as the Proportional
Reporting Ratio and Reporting Odds Ratio, and apply natural language processing
to the free-text narrative fields to extract unstructured clinical detail that
the structured fields do not capture.

% =============================================================================
\section{Thank You \hfill \textcolor{slidenumber}{\small $\sim$15 sec}}
% =============================================================================

Thank you for your time. I am happy to take any questions about the data
pipeline, the statistical methods, or the dashboard implementation. The live
dashboard is deployed on Google Cloud Platform and is accessible at the URL
shown on the title slide if you would like to explore it directly.

% =============================================================================
\vspace{20pt}
\textcolor{accent}{\rule{\linewidth}{1pt}}
\vspace{8pt}

\begin{center}
\begin{tabular}{>{\bfseries}l l}
\toprule
Slide & Estimated Time \\
\midrule
1 -- Title              & 30 sec \\
2 -- What is FAERS?     & 1 min  \\
3 -- Dataset Scope      & 1 min  \\
4 -- Features           & 1 min  \\
5 -- Data Cleaning      & 1.5 min \\
6 -- Feature Engineering & 1 min  \\
7 -- Outlier Detection (methods) & 1 min \\
8 -- Outlier Detection (graphs)  & 30 sec \\
9 -- Normality Testing (methods) & 1 min \\
10 -- Normality Testing (graphs) & 30 sec \\
11 -- PCA               & 1 min  \\
12 -- Pearson Heatmap   & 1 min  \\
13 -- Conclusion        & 1 min  \\
14 -- Thank You         & 15 sec \\
\midrule
\textbf{Total}          & $\approx$ \textbf{12 min} \\
\bottomrule
\end{tabular}
\end{center}

\vspace{6pt}
\begin{notebox}
\textbf{Remaining $\sim$8 minutes:} Use for the live dashboard demo.
Switch to the browser and walk through each tab -- Home, Drug Analysis,
Reaction Explorer, Patient Demographics, Geographic View, Trends \& Timeline,
Severity \& Outcomes, Drug--Reaction Network, Reporter Insights, Safety Signals,
and Statistical Analysis. Let the audience suggest filters to demonstrate
real-time interactivity.
\end{notebox}

\end{document}
